A magnitude do erro cresce linearmente com a leitura $R$ e satisfaz
\[
E(R)=0{,}04R+0{,}8\ \mathrm{V}.
\]
A tensão correta é
\[
T(R)=R-E(R)=0{,}96R-0{,}8.
\]
Isso dá:
\begin{align*}
T(80)&=76{,}0\ \mathrm{V},\\
T(100)&=95{,}2\ \mathrm{V},\\
T(50)&=47{,}2\ \mathrm{V},\\
T(1)&=0{,}16\ \mathrm{V},\\
T(35)&=32{,}8\ \mathrm{V},\\
T(10)&=8{,}8\ \mathrm{V}.
\end{align*}
Os erros percentuais em relação ao valor correto são aproximadamente
\begin{align*}
5{,}3\%,\quad 5{,}0\%,\quad 5{,}9\%,\quad 525\%,\quad 6{,}7\%,\quad 13{,}6\%.
\end{align*}